\documentclass[12 pt, letterpaper]{hmcpset}
\usepackage{lin-alg-preamble}
\begin{document}

\begin{problem}[Dummit \& Foote \S1.7 \#4]
	Let $G$ be a group acting on a set $A$ and fix some $a \in A$. Show that the following sets are subgroups of $G$:
	\begin{enumerate}
		\item the kernel of the action
		\item $\{ g \in G | ga = a \}$ --- this subgroup is called the stabilizer of $a$ in $G$.
	\end{enumerate}
\end{problem}

\begin{solution}
	\begin{enumerate}
		\item
			Call the desired set $\ker(G \to A)$. Suppose we have some arbitrary $g_1, g_2 \in \ker(G \to A)$. We need to show that $g_1 g_2$ and $g_1^{-1}$ are both in $\ker(G \to A)$. Note that $g_1 \cdot a = g_2 \cdot a = a$ for any arbitrary $a \in A$. 

			By computing
			\[
				\begin{split}
					(g_1 g_2) \cdot a & = g_1 \cdot (g_2 \cdot a) \\
							  & = g_1 \cdot a \\
							  & = a
				\end{split}
			\]
			we can see that $g_1 g_2$ fix any arbitrary $a \in A$. That is, $g_1 g_2 \in \ker(G \to A)$. 

			Similarly, by noticing that $a = g_1 \cdot a$ and computing
			\[
				\begin{split}
					g_1^{-1} \cdot a & = g_1^{-1} \cdot (g_1 \cdot a) \\
							 & = (g_1^{-1} g_1) \cdot a \\
							 & = a.
				\end{split}
			\]
			we see that $g_1^{-1}$ fixes any arbitrary $a$. That is, $g_1^{-1} \in \ker(G \to A)$.
		\item
			Call the desired set $\operatorname{stab}_a G$. Suppose we have some arbitrary $g_1, g_2 \in \operatorname{stab}_a G$. We need to show that $g_1 g_2$ and $g_1^{-1}$ are both in $\operatorname{stab}_a G$. Note that $g_1 \cdot a = g_2 \cdot a = a$. 

			By computing
			\[
				\begin{split}
					(g_1 g_2) \cdot a & = g_1 \cdot (g_2 \cdot a) \\
							  & = g_1 \cdot a \\
							  & = a
				\end{split}
			\]
			we can see that $g_1 g_2$ fix $a$. That is, $g_1 g_2 \in \operatorname{stab}_a G$. 

			Similarly, by noticing that $a = g_1 \cdot a$ and computing
			\[
				\begin{split}
					g_1^{-1} \cdot a & = g_1^{-1} \cdot (g_1 \cdot a) \\
							 & = (g_1^{-1} g_1) \cdot a \\
							 & = a.
				\end{split}
			\]
			we get that $g_1^{-1}$ fixes $a$. That is, $g_1^{-1} \in \ker(G \to A)$.
	\end{enumerate}
\end{solution}

\pagebreak

\begin{problem}[Dummit \& Foote \S1.7 \#17]
	Let $G$ be a group and let $G$ act on itself by left conjugation, so each $g \in G$ maps $G$ to $G$ by
	\[
		x \mapsto g x g^{-1}.
	\]
	For fixed $g \in G$, prove that the conjugation by $g$ is an isomorphism from $G$ onto itself (i.e., is an automorphism of $G$). Deduce that $x$ and $g x g^{1}$ have the same order for all $x$ in $G$ and that for any subset $A$ of $G$, $\lvert A \rvert = \lvert g A g^{-1} \rvert$ (here $g A g^{-1} = \{ g a g^{1} | a \in A \} $).
\end{problem}

\begin{solution}
	Define $\sigma_g : G \to G$ by $x \mapsto g x g^{-1}$. First we will show that $\sigma_g$ is a homomorphism. This is just by computing
	\[
		\begin{split}
			\sigma(xy) & = g x y g^{-1} \\
				   & = g x g^{-1} g y g^{-1} \\
				   & = \sigma(x) \sigma(y)
		\end{split}
	\]
	which holds for any $x, y \in G$. 

	Now we will show that $\sigma_g$ is an isomorphism. We will do this by finding an inverse. We claim specifically, that $\sigma_g^{-1}$ is just $\sigma_{g^{-1}}$. Notice that for any $g \in G$
	\[
		\begin{split}
			(\sigma_{g^{-1}} \circ \sigma_g)(x) & = g^{-1} g x g^{-1} g \\
							    & = x.
		\end{split}
	\]
	and since the choice of $g$ was arbitrary and $(g^{-1})^{-1} = g$, we have by the same proof that $(\sigma_g \circ \sigma_{g^{-1}})(x) = x$. Thus, $\sigma_{g^{-1}}$ is $\sigma_g^{-1}$, the unique two sided inverse, of $\sigma_g$, and $\sigma_g$ is an isomorphism. 

	Recall that we have the result that isomorphisms preserve orders. This gives us that $\lvert x \rvert = \lvert \sigma_g(x) \rvert$ and thus, $\lvert x \rvert = \lvert g x g^{-1} \rvert$. 

	Finally, consider the restriction of $\sigma_g$, 
	\[
		{\sigma_g}_{|A} : A \to \sigma_g^\text{img}(A)
	\]
	by $a \mapsto \sigma_g(a)$ and notice that
	\[
		\sigma_g^\text{img}(A) = \{ g a g^{-1} | a \in A \} = g A g^{-1}.
	\]
	By the injectivity of $\sigma_g$ and the fact that we're mapping onto its image, it follows that ${\sigma_g}_{|A}$ is also a bijection, and thus, $\lvert A \rvert = \lvert g A g^{-1} \rvert$.
\end{solution}

\pagebreak

\begin{problem}[Dummit \& Foote \S1.7 \#18]
	Let $H$ be a group acting on a set $A$. Prove that the relation $\sim$ defined by 
	\[
		a \sim b \quad \text{ if and only if } \quad a = h \cdot b \text{ for some } h \in H
	\]
	is an equivalence relation. (For each $x \in A$ the equivalence class of $x$ under $\sim$ is called the \textit{orbit} of $x$ under the action of $H$. The orbits under the action of $H$ partition the set $A$.)
\end{problem}

\begin{solution}
	To show that $\sim$ is an equivalence relation, we need to show reflexivity, symmetry, and transitivity. Let $a, b, c$ be arbitrary elements of $A$, and let $h, h_1, h_2$ be arbitrary elements of $H$.

	We can show reflexivity by noticing that $a = 1_H \cdot a$, and thus $a \sim a$. 

	For symmetry, assume that $a \sim b$. Then we have $a = h \cdot b$. Multiply on both sides by $h^{-1}$ to get $h^{-1} \cdot a = h^{-1} \cdot (h \cdot b)$, which we can rewrite using the associativity of the group action to get $b = h^{-1} \cdot a$, and thus $b \sim a$. 

	Finally, for transitivity, assume that $a \sim b$ and $b \sim c$ such that $a = h_1 \cdot b$ and $b = h_2 \cdot c$, then notice that we have
	\[
		\begin{split}
			a & = h_1 \cdot b \\
			  & = h_1 \cdot (h_2 \cdot c) \\
			  & = (h_1 h_2) \cdot c
		\end{split}
	\]
	and thus, we have $a \sim c$. 

	Since we have shown, reflexivity, symmetry, and transitivity, we have that $\sim$ is an equivalence relation.
\end{solution}

\pagebreak

\begin{problem}[Dummit \& Foote \S1.7 \#19]
	Let $H$ be a subgroup of the finite group $G$ and let $G$ act on $G$ (here $A = G$) by left multiplication. Let $x \in G$ and let $\mathcal{O}_H^x$ be the orbit of $x$ under the action of $H$. Prove that the map $H \to \mathcal{O}_H^x$ defined by $h \mapsto hx$ is a bijection (hence all orbits have cardinality $\lvert H \rvert$). From this and the preceding exercise deduce \textit{Lagrange's Theorem}:
	\begin{center}
		\textit{if $G$ is a finite group and $H$ is a subgroup of $G$ then $\lvert H \rvert$ divides $\lvert G \rvert$.}
	\end{center}
\end{problem}

\begin{solution}
	Let $\varphi$ denote the map $H \to \mathcal{O}_H^x$ by $h \mapsto hx$. First we will show that $\varphi$ is surjective, and then that is injective. 

	Notice that by definition, the orbit of $x$ under $H$, $\mathcal{O}_H^x$ is all elements of the form $hx$ since, for any $g \in G$, $x \sim g$ if and only if $g = h x$. Since $\varphi^\text{img}$ is just all of these elements of the form $hx$, we have
	\[
		\begin{split}
			\varphi^\text{img}(H) & = \{ hx | h \in H \} \\
					      & = \mathcal{O}_H^x.
		\end{split}
	\]
	Thus, $\varphi$ is surjective. 

	To show injectivity, we suppose $\varphi(h_1) = \varphi(h_2)$ for some $h_1, h_2 \in H$. Then we have $h_1 x = h_2 x$. Right cancellation of $x$ then gives us $h_1 = h_2$. Thus, $\varphi$ is injective. 

	Since $\varphi$ is injective and surjective, it is bijective, and thus, $\lvert H \rvert = \lvert \mathcal{O}_H^x \rvert$, for any $x \in G$. That is, the order of the orbit of any element under the action of $H$ is just the order of $H$. 

	Notice that this also means that we have $\lvert \mathcal{O}_H^x \rvert = \lvert \mathcal{O}_H^y \rvert$ for any $x, y \in G$. The orders of the orbits of any two elements are the same --- they're both the order of $H$. By the previous exercise we have that the orbits of elements of $G$ under the action of $H$ partition the group. That is, we have some $x_1, \dots, x_m \in G$, all in different orbits under the action of $H$ such that
	\[
		\coprod_{i = 1}^m \mathcal{O}_H^{x_i} = G 
	\]
	(their disjoint union is $G$), and thus,
	\[
		\sum_{i = 1}^m \lvert \mathcal{O}_H^{x_i} \rvert = \lvert G \rvert.
	\] 

	Since we have $\lvert \mathcal{O}_H^{x_i} \rvert = \lvert H \rvert$ this is just equivalent to $m \lvert H \rvert = G$ and thus, $\lvert H \rvert$ divides $\lvert G \rvert$.
\end{solution}

\pagebreak

\end{document}
